\small\begin{longtable}{p{0.07\textwidth}p{0.08\textwidth}p{0.72\textwidth}}
\caption{Parecer individual sobre as 66 indicações. A: representação autenticada; L: referência lexical; B: identificação bibliográfica; I: exemplar específico; S: sem ocorrência autenticada.}\\
\toprule ID & Estado & Resultado e limites \\\midrule\endfirsthead
\toprule ID & Estado & Resultado e limites \\\midrule\endhead
R01 & I & Exemplar de 25/07/1904, n. 207, recuperado; busca no OCR das quatro páginas não autenticou a frase do ovo. OCR degradado impede conclusão negativa definitiva.\\
R02 & S & Busca por título, década e riso não individualizou artigo, edição e página. A descrição genérica não comprova grafia.\\
R03 & S & Sem ano e edição do almanaque ou título da anedota. Busca não individualizou ocorrência; é necessário delimitar volume e seção.\\
R04 & L & Texto integral recuperado: expressão de riso, PDF p. 28. Referência lexical; não autentica o diálogo atribuído ao colono Maciel.\\
R05 & L & Exemplar consultado: segunda edição, 1917, PDF p. 26: Riu-se muito. Não datar a passagem de 1909 sem cotejo da primeira edição.\\
R06 & L & Oito capítulos consultados; capítulo VI contém risos casquinados. É descrição lexical, não uma sequência onomatopaica demonstrada.\\
R07 & S & A busca recuperou Quá em Quase doutor, Careta, 08/05/1915, via Vida urbana. Contexto de recusa, não riso: falso positivo excluído.\\
R08 & I & Edição n. 526 de 12/10/1912 confirmada na capa. Busca nas 58 páginas do OCR não autenticou prefeito/carne barata; desenhos exigem revisão visual integral.\\
R09 & S & Pista alternativa: poema Epístola, A Cigarra, n. 16, 20/01/1915, p. 27, em reprodução externa. Não conferida no exemplar original; não incluída no corpus.\\
R10 & L & Texto de 1915 consultado. Terceira parte, cap. II, contém Você, Quaresma, é um visionario; sem Ah! ah! e sem a continuação sobre plantar. Citação recebida não é transcrição fiel.\\
R11 & A & Quinta edição de 1919: p. impressa 19, PDF p. 24: Quá, quá, quá!, com risos explícitos no contexto. Não retrodatado automaticamente a 1918.\\
R12 & L & Monólogo de uma sombra: Chorando e rindo na ironia infausta. Confirma referência temática, não onomatopeia.\\
R13 & S & Frase de 14/02/1925 não autenticada. Fon-Fon!, maio de 1925, p. 3, noticia retorno do Correio após suspensão; é indício de problema cronológico, não prova autônoma de que não existiu edição em fevereiro.\\
R14 & A & 17/06/1928, n. 290, p. 4: (risos, muito bem), em discurso de Mattos Pimenta. Rubrica descritiva, não grafia de gargalhada.\\
R15 & S & Busca por título, década e risos não individualizou reportagem verificável; faltam dia e página.\\
R16 & A & Edição UFFS de 2019, p. 139: Quá! quá! quá quaquá!... e indicação de gargalhada. Obra original de 1928; a passagem não é a frase Ah! ah! enviada.\\
R17 & L & Exemplar de 1924, p. 114, PDF p. 116: morreu nos lábios de todos o sorriso. Referência lexical; não comprova onomatopeia.\\
R18 & B & BBM-USP registra edição de 1954. A alegação de composição nos anos 1920 não foi demonstrada. Nenhuma passagem de riso autenticada.\\
R19 & I & PDF identificado pelo acervo como n. 23 de 1935 recuperado; capa e OCR consultados. Data de 08/06 e frase sobre rei da Inglaterra não autenticadas; falta conferir expediente legível.\\
R20 & S & Exemplar n. 1175 de 1935 recuperado e OCR pesquisado. Há candidatos lexicais, mas nenhuma grafia atribuída autenticada por imagem.\\
R21 & S & Exemplar n. 31 de 1930 recuperado e OCR pesquisado. Referências lexicais candidatas não autenticam a generalização sobre marcações de riso.\\
R22 & B & Editora Record confirma estreia em 1930; aquisição automática recebeu HTTP 429. Não se confirmou passagem nem grafia.\\
R23 & L & Provas da segunda edição, 1938, recuperadas; PDF p. 63: sorrisos frios. A frase latim de missa não foi localizada no OCR; não retrodatado a 1934.\\
R24 & S & Reforma Agrária (1934), de Marquie: buscas por título, ano e variante do nome não identificaram registro bibliográfico suficiente. Não tratar como obra inexistente.\\
R25 & I & Edição de 18/03/1944, n. 6564, confirmada no cabeçalho. Busca do OCR não autenticou ministro/boatos de racionamento; texto de baixa qualidade.\\
R26 & S & Busca por título, ano e rubrica não individualizou transmissão ou reportagem com página. Não demonstrada a descrição recebida.\\
R27 & S & Busca produziu falso positivo: Jornal de Notícias, de 02/03/1947, e não Folha da Manhã. Documento incluído separadamente; não substitui esta referência.\\
R28 & B & Bibliografia da Academia Brasileira de Letras confirma Fogo morto, 1943. Texto e passagem sobre riso não conferidos.\\
R29 & B & Editora identifica O Continente na trilogia publicada entre 1949 e 1962. Metadado não autentica diálogo de caserna ou taberna.\\
R30 & B & Museu Casa Guimarães Rosa registra exemplar de Sagarana de 1946. Frase He! he! e mundo dá muitas voltas não autenticada.\\
R31 & S & Revista de 1950 e alegado Almanaque de 1953 não são a mesma referência. Faltam número, história e página para verificar moedas e Quá-quá-quá.\\
R32 & S & Busca de frase e jornal não localizou plano é honesto. Ano de 1955 sem dia, coluna e página não individualiza o documento.\\
R33 & S & Busca por jornal, década e risos não individualizou crônica. É necessária referência de edição.\\
R34 & B & Página editorial de Grande sertão: veredas consultada; primeira edição em 1956. Nenhuma passagem específica de riso foi autenticada.\\
R35 & B & Registro institucional sobre Vestido de noiva identifica escrita em 1943; estreia e publicação são eventos distintos. A edição em livro dos anos 1950 não foi individualizada.\\
R36 & S & Crônicas de Stanislaw Ponte Preta não designa título, edição ou crônica específica. Identidade do cronista não confirma data e grafia de passagem.\\
R37 & S & Coleção identificada no dossiê da Biblioteca Nacional. URL tentada para n. 12 retornou 404; diálogo atribuído a Jaguar permanece não autenticado.\\
R38 & S & A reforma gráfica do jornal não demonstra uso de hahaha. Busca não localizou entrevista e página específicas na década de 1960.\\
R39 & S & Sem edição e entrevistado identificados. Busca por título, década e rubrica não autenticou onomatopeia corrida.\\
R40 & B & Página editorial de Quarup localizada; download automático HTTP 429. Frase sobre jornais não autenticada e sem página.\\
R41 & B & Registro museológico localizado para A paixão segundo G.H.; download retorna página sem conteúdo textual suficiente. Não confirma sorriso mental ou grafia.\\
R42 & B & ABL distingue edição italiana de 1974 e portuguesa de 1975. Alegada composição nos anos 1960 não autoriza contar publicação naquela década.\\
R43 & S & Menção genérica a gibis dos anos 1970. Busca não individualizou número e história com as formas atribuídas.\\
R44 & S & Busca por Opinião, década e risos não individualizou passagem verificável. A caracterização ideológica não prova grafia.\\
R45 & S & Busca por Movimento, década e risos não individualizou debate com localizador. Pendente identificação de edição.\\
R46 & B & Página editorial de As meninas consultada. Frase sobre aposta e Ha-ha-ha! não autenticada; sem acesso integral conferido.\\
R47 & B & Página editorial localizada na busca; aquisição automática falhou com 404. Data de obra não é prova de uma risada específica.\\
R48 & B & Coletânea de 1975 identificada também em pesquisa institucional UFRN; catálogo WorldCat bloqueou aquisição com 403. Nenhuma passagem de riso confirmada.\\
R49 & S & Busca da data 14/11/1986 e frase de congelamento não autenticou Rs. na seção Painel. Rubrica risos de outra publicação de 1994 não substitui esta alegação.\\
R50 & S & Busca por Zózimo, década e rs não individualizou nota. Não demonstrado uso comum de rs no jornal nos anos 1980.\\
R51 & S & Busca por jornal, década e rsrs não individualizou nota. Alegada economia da datilografia permanece explicação sem evidência específica.\\
R52 & B & Página editorial de Viva o povo brasileiro consultada. Descrição literária genérica não comprova forma particular de riso.\\
R53 & B & Página editorial de O Alquimista consultada. Sorriso e rir são itens lexicais; não devem ser classificados automaticamente como interjeições.\\
R54 & B & Página editorial de Feliz ano velho consultada. Frase burradas do colégio e sequência Ha ha ha não autenticadas.\\
R55 & S & Busca de Veja, 07/05/1993 e frase não fugirei da raia não individualizou entrevista. Data e entrevistado continuam sem autenticação.\\
R56 & S & Busca de Folhateen, junho de 1997, mIRC e grafias não identificou reportagem com dia e página. Não confirmada a frase recebida.\\
R57 & S & Busca de Bizz, década e hahaha não identificou entrevista específica. Sem edição, autoria e página.\\
R58 & B & Fonte institucional UEPA registra primeira publicação em 1941. Ano de 1998 pode designar edição; não é data de nascimento da obra nem da passagem.\\
R59 & B & Página editorial informa lançamento em 10/09/2003. Planejamento nos anos 1990 não demonstrado e não equivale a publicação.\\
R60 & B & Página editorial informa lançamento em 02/08/1991. Nenhuma passagem sobre riso foi conferida no texto integral.\\
R61 & S & Busca de Capricho, setembro de 2007 e frase quase quebrou a cara não individualizou reportagem. Não se sustenta afirmação de primeira impressão de kkkk.\\
R62 & S & Caderno de Informática de dois jornais é referência ambígua. Busca por internetês e kkkk não individualizou tabela de abreviações com data e página.\\
R63 & S & Busca por Meia Hora, 2005 e kkkk não identificou manchete. Fundação do veículo não demonstra adoção de grafia específica.\\
R64 & S & Título, autor e editora do alegado livro de 2005 não identificados. Resultado homônimo é postagem de blog de 2008, que não autentica o livro.\\
R65 & S & Manual do MSN Messenger sem autor, ISBN ou editora. Busca não individualizou obra e edição; não autentica convenções descritas.\\
R66 & S & Literatura de Blog é categoria, não referência bibliográfica individual. Cardosonline e Melhores do Mundo não podem ser convertidos em livro sem título, autor e edição.\\
\bottomrule\end{longtable}\normalsize